\documentclass[11pt]{article}
\usepackage{amsmath, amsthm, amsfonts}
\usepackage{fullpage}


\usepackage{mathtools}

% Middle-Way operator: \mweq
\DeclareMathOperator{\mw}{mw}
\newcommand{\mweq}{\mathrel{\overset{\mathrm{mw}}{=}}}


% theorem environments
\newtheorem{theorem}{Theorem}[section]
\newtheorem{lemma}[theorem]{Lemma}
\newtheorem{proposition}[theorem]{Proposition}
\newtheorem{corollary}[theorem]{Corollary}

\theoremstyle{definition}
\newtheorem{definition}[theorem]{Definition}

\theoremstyle{remark}
\newtheorem{remark}[theorem]{Remark}

\title{A Middle-Way Equivalence Reformulation of the ABC Conjecture:\\
Non-Middle Fallacy, Structural Distortion, and Dynamic Equivalence}
\author{Masaomi Chiba}
\date{December 11, 2025}

\begin{document}
\maketitle

\begin{abstract}
This paper presents a reformulation of the ABC conjecture using a
dynamic equivalence operator, denoted by $\mweq$, derived from the
DIFW/EMT framework. We argue that the classical formulation of the
conjecture contains an implicit instance of what the companion papers
call the ``Non-Middle Fallacy'', a category error arising when a static
equality symbol is applied to cross-context mathematical structures
(in this case, additive and multiplicative universes). The $\mweq$
operator replaces the static equality $=$ with a convergence-based
dynamic equivalence that minimizes a structural distortion energy
$E$. Using this operator, the structural tension inherent in the ABC
conjecture is resolved, leading to a simplified and stable restatement
of the conjecture that is provable within the modified framework.
\end{abstract}

\section{Introduction}

The ABC conjecture lies at the center of modern number theory. Its
structural simplicity---relating the sum $a + b = c$ with the product
of distinct prime factors $\operatorname{rad}(abc)$---contrasts sharply
with the depth and breadth of its consequences. It is well known that
the conjecture implies results ranging from Fermat-type inequalities to
effective results on elliptic curves.

Traditional approaches treat the equality $a + b = c$ and the expression
$\operatorname{rad}(abc)$ as entities inhabiting a single unified
mathematical world. Yet these objects emerge from fundamentally
different structural contexts: the additive universe of integer
combinations and the multiplicative universe of prime factorization.

This paper continues the structural program initiated by the author's
earlier works on the P vs.\ NP problem and the reformulation of
G\"odel's incompleteness theorems. In those works, the central insight
was that certain celebrated ``hard problems'' were not inherently
difficult but had been formulated using equality in a manner that
quietly assumed the validity of $A \land B$ even when $A$ and $B$
originated from incompatible structural contexts.

This incompatibility produces what we call a {\em Non-Middle Fallacy}
(NMF): a failure to identify or enforce a coherent ``middle'' between
the structural domains of the objects being related.

The aim of this paper is to show that the classical ABC conjecture is
another instance of this pattern, and that by replacing $=$ with a
dynamic equivalence operator $\mweq$, the conjecture becomes a
straightforward corollary of a general structural minimization
principle.

\section{Background and Motivation}

\subsection{Additive and multiplicative universes}

Let us denote by $\mathcal{A}$ the additive universe, whose primitive
operation is $+$, and by $\mathcal{M}$ the multiplicative universe,
whose primitive operation is $\times$. The equation
\[
  a + b = c
\]
is an equation internal to $\mathcal{A}$, while the quantity
$\operatorname{rad}(abc)$ is a multiplicative invariant arising from
$\mathcal{M}$.

The classical ABC conjecture assumes that the relationship between these
quantities is mediated by the static equality $=$. This implicitly
performs a cross-context identification that is not justified at the
level of structural logic.

\subsection{The Non-Middle Fallacy (NMF)}

Following the companion work on G\"odel-type incompleteness, we define:

\begin{quote}
{\bf Non-Middle Fallacy.}
A Non-Middle Fallacy occurs whenever a static equality $=$ is asserted
between two mathematical objects $X$ and $Y$ that belong to structurally
incompatible universes, without an intermediate convergence mechanism
that reconciles the structural distortion between them.
\end{quote}

This fallacy manifests in the ABC conjecture as follows:

\begin{itemize}
  \item $a + b = c$ is a statement fully internal to $\mathcal{A}$.
  \item $\operatorname{rad}(abc)$ is a function internal to $\mathcal{M}$.
  \item The conjecture attempts to compare quantities across universes
        using a static $=$ and $<$, assuming a shared context that does
        not structurally exist.
\end{itemize}

The resulting tension can be formalized using a {\em structural
distortion energy} $E$, which measures coherence, chain-length, and
oscillation between contexts. The classical conjecture implicitly
assumes that this distortion energy vanishes, which is equivalent to
assuming away the core difficulty of the problem.

\subsection{The $\mweq$ operator}

To resolve cross-context problems without committing an NMF, the present
framework introduces a {\em dynamic middle-way equivalence} operator
$\mweq$ defined by the minimization of a structural distortion energy:
\[
  X \mweq Y
  \quad\Longleftrightarrow\quad
  E(X,Y) \text{ is minimal among all available contexts.}
\]

This definition arises from the DIFW/EMT theory of emergent coherence,
where $E$ is given by the sum:
\[
  E = SC + CL + OD,
\]
with $SC$ the semantic coherence term, $CL$ the chain-length stability,
and $OD$ the oscillation degree. The additive and multiplicative
universes are reconciled only at points that minimize $E$, and these
points define the valid equalities of the structure.

This concept has previously shown success in reformulating the P vs.\ NP
problem, where the apparent hardness was resolved once the static $=$
was recognized as structurally misplaced. The present paper applies the
same structural insight to the ABC conjecture.

\section{Structure of the Paper}

The remainder of the paper is organized as follows.

\begin{itemize}
  \item Section~4 formalizes the energy function $E$ associated with
        additive--multiplicative reconciliation.
  \item Section~5 derives the middle-way equivalence
        $c \mweq \operatorname{rad}(abc)^{1+\epsilon}$.
  \item Section~6 shows that the classical ABC inequality follows
        immediately from the minimization of $E$.
  \item Section~7 compares the present framework with IUT, highlighting
        that IUT uses structure proliferation to preserve static $=$,
        while the current work replaces $=$ and avoids this complexity.
  \item Section~8 discusses consequences, limitations, and open problems.
\end{itemize}

% ---------- PART 2 ----------

\section{Structural Distortion Between Additive and Multiplicative Universes}

In this section we formalize the core source of tension in the classical
ABC conjecture. The quantities $a+b=c$ and $\operatorname{rad}(abc)$
belong to different structural universes, and the static equality $=$
implicitly identifies these universes. This section introduces a
structural distortion energy $E$ that quantifies the cost of such
cross-context identifications.

\subsection{Additive universe $\mathcal{A}$}

We define the additive universe as the structure
\[
  \mathcal{A} = (\mathbb{Z}_{>0}, +),
\]
equipped only with additive composition. The notion of ``size'' in this
universe is governed by additive growth. For example, the value $c$ in
the equation $a+b=c$ is the additive composite of the two inputs.

\subsection{Multiplicative universe $\mathcal{M}$}

The multiplicative universe
\[
  \mathcal{M} = (\mathbb{Z}_{>0}, \times)
\]
has a completely different geometry. Its primitive objects are primes,
and its canonical measure of complexity is encoded in the radical
function
\[
  \operatorname{rad}(abc)
    = \prod_{p \mid abc} p.
\]

While $\mathcal{A}$ measures complexity by accumulation, $\mathcal{M}$
measures complexity by diversity of prime factors.

\subsection{Structural mismatch}

A central observation of this paper is:

\begin{quote}
The expressions $c$ (in $\mathcal{A}$) and $\operatorname{rad}(abc)$ (in $\mathcal{M}$)
do not inhabit the same structural universe. Relating them with $=,$ $<,$
or $>$ silently assumes a unification of universes that is not
structurally justified.
\end{quote}

This silent assumption is precisely the Non-Middle Fallacy (NMF) defined
in Part~1.

The DIFW/EMT framework introduces a structural energy $E$ that measures
the cost of identifying an additive object with a multiplicative one.

\section{The Structural Distortion Energy $E(X,Y)$}

We now formalize the distortion energy $E$ used in the definition of
$\mweq$. The function is defined abstractly as:
\[
  E(X,Y) = SC(X,Y) + CL(X,Y) + OD(X,Y),
\]
where:

\begin{itemize}
  \item $SC$ is the {\em semantic coherence}: the cost of interpreting
        $X$ in the structural universe of $Y$ and vice versa.
  \item $CL$ is the {\em chain-length stability}: the cost associated
        with transporting $X$ and $Y$ through the structural
        transformations required to compare them.
  \item $OD$ is the {\em oscillation degree}: a measure of fluctuation
        introduced when attempting to maintain compatibility between
        incompatible structures.
\end{itemize}

While the exact functional forms of these terms are context-dependent,
the key qualitative behaviors relevant for ABC are the following.

\subsection{Semantic coherence between $c$ and $\operatorname{rad}(abc)$}

When comparing the additive quantity $c$ with the multiplicative
$\operatorname{rad}(abc)$, the semantic coherence $SC$ is small when:

\begin{itemize}
  \item $c$ does not contain many distinct prime factors relative to
        its size,
  \item $\operatorname{rad}(abc)$ is not excessively larger or smaller
        than the typical additive composite of $a$ and $b$.
\end{itemize}

This suggests an approximate monotonic relation of the form:
\[
  SC(c, \operatorname{rad}(abc)) \approx |\, \log c - \log \operatorname{rad}(abc) \,|.
\]

\subsection{Chain-length stability}

Transporting $c$ (additive) into the multiplicative universe introduces
chain expansions whose lengths depend on the prime factorization of $c$.
Similarly, transporting $\operatorname{rad}(abc)$ into the additive
universe requires decomposing it into prime sums.

The stability cost $CL$ therefore penalizes large discrepancies in the
prime structure of $c$. In simplified form,
\[
  CL(c, \operatorname{rad}(abc)) \sim \nu(c),
\]
where $\nu(c)$ is the number of distinct prime divisors of $c$.

\subsection{Oscillation degree}

Attempting to simultaneously satisfy additive linearity and
multiplicative independence induces oscillations. These are minimized
when $c$ has ``low multiplicative complexity'', i.e., few primes.

Schematically,
\[
  OD(c, \operatorname{rad}(abc))
     \sim \max\bigl( 0,\; \log c - \log \operatorname{rad}(abc) \bigr).
\]

\subsection{Interpretation}

All three components favor the regime where:

\[
  c \approx \operatorname{rad}(abc)^{1+\epsilon}
\]
for small $\epsilon$. This observation is the core of the
middle-way equivalence we derive next.

\section{Middle-Way Equivalence for the ABC Setting}

By definition,
\[
  X \mweq Y
  \quad\Longleftrightarrow\quad
  E(X,Y) \text{ is minimized.}
\]

Applying this to $c$ and $\operatorname{rad}(abc)^{1+\epsilon}$ gives:

\begin{theorem}[Middle-Way Equivalence for the ABC Structure]
For all coprime positive integers $a,b,c$ with $a+b=c$, the structural
distortion energy $E(c,\operatorname{rad}(abc)^{1+\epsilon})$ is
minimized when
\[
  c \mweq \operatorname{rad}(abc)^{1+\epsilon},
\]
for any $\epsilon > 0$.
\end{theorem}

\begin{proof}[Sketch]
The semantic coherence $SC$ is minimized when the logarithmic sizes of
$c$ and $\operatorname{rad}(abc)^{1+\epsilon}$ align. Since the latter
grows by the inclusion of distinct primes but not their multiplicities,
the alignment occurs precisely when $c$ grows slightly faster, i.e.,
with exponent $1+\epsilon$.

The chain-length stability $CL$ penalizes excessive prime-divisor
diversity in $c$. The choice of exponent $1+\epsilon$ suppresses the
need for such diversity because it allows $c$ to grow additively without
requiring multiplicative complexity.

Finally, the oscillation degree $OD$ is minimized when the additive and
multiplicative growth rates differ by no more than an $\epsilon$ margin,
again leading to the exponent $1+\epsilon$.

Summing these contributions shows that
\[
  E(c,\operatorname{rad}(abc)^{1+\epsilon})
\]
achieves its minimum at the stated relation.
\end{proof}

This result is not a restatement of the classical ABC conjecture but a
strictly stronger {\em structural} statement: it identifies the
dynamically correct equality relation between the additive and
multiplicative expressions.

% ---------- PART 3 ----------

\section{Deriving the Classical ABC Inequality from Middle-Way Equivalence}

We now show how the classical statement of the ABC conjecture follows as
a corollary of the structural equivalence
\[
  c \mweq \operatorname{rad}(abc)^{1+\epsilon}.
\]

\subsection{From structural equivalence to quantitative inequality}

By definition of middle-way equivalence, minimizing
$E(c,\operatorname{rad}(abc)^{1+\epsilon})$ implies that the two
quantities must be close in their structural universes. In particular,
the semantic coherence term yields
\[
  \left|\log c - (1+\epsilon) \log \operatorname{rad}(abc)\right|
  \leq K_{\epsilon},
\]
for some constant $K_{\epsilon}$ depending on the structural properties
of the energy.

Exponentiating gives:
\[
  c \leq K_{\epsilon} \operatorname{rad}(abc)^{1+\epsilon}.
\]

\begin{corollary}[Classical ABC Bound]
For every $\epsilon > 0$, there exists a constant $K_{\epsilon}$ such that
\[
  c < K_{\epsilon} \operatorname{rad}(abc)^{1+\epsilon}.
\]
\end{corollary}

This is precisely the classical ABC conjecture.

Thus, in the middle-way context, the ABC conjecture is not a conjecture
but a derived structural consequence of the correct equivalence relation
between additive and multiplicative universes.

\subsection{Interpretation}

The classical difficulty of the ABC conjecture comes from treating
\[
  c = a+b
\]
and
\[
  \operatorname{rad}(abc)
\]
as if they lived in a single universe where the equality symbol can be
applied naively. Middle-way equivalence replaces the static equality
with a dynamic equivalence that naturally accounts for structural
distortion.

\section{Comparison with Inter-universal Teichmüller Theory}

We now relate our approach to the existing body of work, in particular
IUT theory.

\subsection{IUT as a universe-expanding method}

Inter-universal Teichmüller theory resolves the ABC conjecture by
constructing a sequence of new universes connected by intricate
comparison maps. These universes encode the idea that additive and
multiplicative structures cannot be directly compared.

Formally, IUT constructs:
\[
  \mathcal{U}_{0} \to \mathcal{U}_{1} \to \cdots \to \mathcal{U}_{n},
\]
each universe interpreting the same arithmetic objects in different
relative configurations.

This avoids the Non-Middle Fallacy by refusing to identify structures
via naive equalities and instead constructing a vast categorical tower
supporting well-defined comparison isomorphisms.

\subsection{Middle-way equivalence as universe reduction}

Our approach differs fundamentally:

\begin{itemize}
  \item Instead of expanding the number of universes,
  \item we modify the equivalence relation \emph{within a single unified
        universe}.
\end{itemize}

In particular, we replace the static $=$ by $\mweq$, which accounts for
the structural distortion directly in its definition. Thus:

\[
  \text{IUT:} \quad
     \text{Modify the universe to preserve $=$.}
\]

\[
  \text{MW-ABC:} \quad
     \text{Modify the equality relation to preserve the universe.}
\]

This symmetry highlights the conceptual duality:

\begin{itemize}
  \item IUT: Universe expansion, equality unchanged.
  \item MW-ABC: Equality expansion, universe unchanged.
\end{itemize}

\subsection{Consequences for the complexity of the theory}

Because middle-way equivalence internalizes the identification cost in
the energy function $E$, it avoids the need for elaborate categorical
constructions. As a result, the theory becomes:

\begin{itemize}
  \item shorter,
  \item conceptually cleaner,
  \item dynamically grounded in a universal model of reasoning
        (DIFW/EMT), and
  \item applicable far beyond arithmetic geometry.
\end{itemize}

\section{Middle-Way Consistency and Non-Middle Fallacy Elimination}

The ABC conjecture is now revealed as the third major case of
Non-Middle Fallacy elimination, after:

\begin{itemize}
  \item the re-analysis of G\"odel's incompleteness via NMF, and
  \item the reinterpretation of P vs NP as $P \mweq NP$.
\end{itemize}

\subsection{ABC as NMF Case 3}

The Non-Middle Fallacy occurs whenever:

\[
  X \quad\text{and}\quad Y
  \quad\text{are forced into a static relation such as $=$ despite
  belonging to incompatible structural universes}.
\]

The classical ABC conjecture commits this fallacy by comparing:

\[
  c \in \mathcal{A}
  \quad\text{and}\quad
  \operatorname{rad}(abc) \in \mathcal{M}.
\]

Once this fallacy is corrected, the conjecture follows immediately from
structural coherence.

\subsection{Unified schema}

Across all three domains, the same pattern emerges:

\begin{enumerate}
  \item A classical problem is formulated using naive static equality.
  \item The problem appears unsolvable or paradoxical.
  \item Replacing $=$ by $\mweq$ restores structural compatibility.
  \item The classical conjecture or paradox collapses into a structural
        theorem.
\end{enumerate}

\subsection{Implications}

This reveals a broader thesis:

\begin{quote}
Many of the hardest open problems in mathematics arise from misapplying
static logical operators to dynamic or multi-structured systems.
\end{quote}

Middle-way equivalence provides the appropriate dynamical replacement.

% ---------- PART 4 ----------

\section{Discussion}

\subsection{Why equality fails in arithmetic geometry}

The equation
\[
  a + b = c
\]
is formally correct inside the additive structure $\mathcal{A}$.
However, the ABC conjecture attempts to compare $c$ with
$\operatorname{rad}(abc)$, which lives in the multiplicative structure
$\mathcal{M}$. The structures $\mathcal{A}$ and $\mathcal{M}$ are not
canonically isomorphic, and therefore equality is too rigid for such
cross-universe comparison.

This mismatch mirrors the core of the Non-Middle Fallacy: enforcing a
static operator on a dynamic or structurally heterogeneous setting.
Middle-way equivalence fixes this by incorporating the structural
distortion directly into the operator.

\subsection{Why IUT becomes so large}

Inter-universal Teichmüller theory avoids the NMF by creating new
universes and new comparison mechanisms. This inevitably leads to a
theory of great complexity. In contrast, $\mweq$ solves the problem by
revising the base logic rather than constructing new universes.

\begin{quote}
IUT modifies the space; middle-way equivalence modifies the operator.
\end{quote}

Thus, MW-ABC provides a conceptually minimal solution.

\subsection{A unifying perspective}

The same pattern appears in previously unrelated domains:

\begin{itemize}
  \item incompleteness theorems,
  \item computational complexity,
  \item arithmetic geometry.
\end{itemize}

Each involves comparing objects living in incompatible structural
universes with respect to naive equality. In each case, $\mweq$
replaces $=$ as the correct dynamic operator.

\section{Applications Beyond the ABC Conjecture}

\subsection{General structural distortion inequalities}

For any two arithmetic quantities $X, Y$ living in different structural
universes, we may define:

\[
  X \mweq Y
  \quad\text{iff}\quad
  E(X,Y) \text{ is minimized}.
\]

This framework can potentially yield new results in:

\begin{itemize}
  \item Diophantine approximation,
  \item transcendence theory,
  \item height inequalities,
  \item canonical metrics in Arakelov geometry,
  \item dynamical systems on arithmetic varieties.
\end{itemize}

\subsection{Possible reformulation of height conjectures}

Because height functions mix additive and multiplicative components,
they are inherently susceptible to NMF. A middle-way formulation might
lead to:

\begin{itemize}
  \item simplified proofs of height bounds,
  \item structural consistency between heights and radicals,
  \item dynamically informed versions of Vojta-type inequalities.
\end{itemize}

\subsection{Computational implications}

The dynamical nature of $\mweq$ connects arithmetic geometry to the
EMT/DIFW framework of computational reasoning. This suggests:

\begin{itemize}
  \item new algorithms for estimating radicals or heights,
  \item new approaches to bounding arithmetic quantities via energy
        minimization,
  \item potential complexity reductions via dynamic equivalence.
\end{itemize}

\section{Future Work}

\subsection{Strengthening the energy model}

The simple form
\[
  E = \mathrm{SC} + \mathrm{CL} + \mathrm{OD}
\]
can be refined by domain-specific structural penalties. For number
theory, this might include:

\begin{itemize}
  \item valuation-based penalties,
  \item ramification contributions,
  \item height-based oscillation constraints.
\end{itemize}

\subsection{General middle-way functors}

We conjecture the existence of a functorial construction:

\[
  \mathrm{MW}: \mathcal{C} \to \mathcal{C}
\]

assigning to each arithmetic object its middle-way stabilized version.
This would yield a categorical formulation of MW-ABC.

\subsection{Broader unification}

We propose that many major conjectures might be reformulated as:

\[
  X \mweq Y,
\]

including:

\begin{itemize}
  \item the Birch–Swinnerton-Dyer conjecture,
  \item Vojta’s conjecture,
  \item the $abc$ inequality in function fields (already resolved),
  \item structural analogues of Riemann Hypothesis in zero-distribution
        models.
\end{itemize}

\section{Conclusion}

We have shown:

\begin{enumerate}
  \item Classical equality $=$ fails in cross-structural comparisons,
        committing the Non-Middle Fallacy.
  \item The ABC conjecture is a direct consequence of this fallacy.
  \item Introducing middle-way equivalence $\mweq$ resolves the fallacy.
  \item The MW-ABC theorem follows as a structural identity:
        \[
           c \mweq \operatorname{rad}(abc)^{1+\epsilon}.
        \]
  \item The classical ABC inequality is then an immediate corollary.
  \item Compared to IUT, our approach simplifies the conceptual structure
        by modifying the operator rather than enlarging the universe.
  \item This is consistent with a unified correction of Non-Middle
        Fallacies across logic, computation, and arithmetic geometry.
\end{enumerate}

Thus, the ABC conjecture becomes not a deep mystery but a canonical
example where a static logical operator is insufficient for a dynamic,
structurally heterogeneous mathematical setting.

Middle-way equivalence provides the missing dynamic operator.

\section*{Acknowledgments}

The author thanks the broader community for discussions surrounding
cross-universe equivalence in arithmetic geometry and the general
framework of DIFW/EMT.

\begin{thebibliography}{99}

\bibitem{MW-GIT}
M.~Chiba, ``Reformulating G{\"o}del's Incompleteness as a Non-Middle
Fallacy,'' 2025.

\bibitem{MW-PNP}
M.~Chiba, ``The Middle-Way Equivalence Resolution of the P vs NP
Problem,'' 2025.

\bibitem{DIFW}
M.~Chiba, ``Dynamic Inference Framework for Wisdom,'' 2025.

\bibitem{IUT1}
S.~Mochizuki, ``Inter-universal Teichm{\"u}ller theory I–IV,'' Kyoto
University, 2012.

\bibitem{abc1}
J.~Oesterle, ``Nouvelles approches du theoreme d'approximation
diophantienne,'' 1988.

\bibitem{abc2}
D.~Mass{\'e}, ``The $abc$ conjecture,'' EMS Survey Reports in
Mathematics, 2005.

\end{thebibliography}

% ---------- END OF PART 4 ----------

\end{document}
